\documentclass[conference]{IEEEtran}
\IEEEoverridecommandlockouts

\usepackage{cite}
\usepackage{amsmath,amssymb,amsfonts}
\usepackage{graphicx}
\usepackage{xspace}
\usepackage{complexity}
\usepackage{color}
\usepackage{comment}
\usepackage{multirow}

\graphicspath{{figs/}}
\newcommand{\todo}[1]{\textcolor{red}{[TODO: #1]}}
\newenvironment{listing}[1]{%
        \begin{list}{*}{%
%                \let\makelabel\namelistlabel
                 \settowidth{\labelwidth}{#1}%
                 \setlength{\leftmargin}{\labelwidth}%
                 \advance \leftmargin by 12pt
%                  \setlength{\labelsep}{5pt}%
                   \setlength{\itemsep}{0pt}%
%                   \setlength{\parsep}{0pt}%
%                   \setlength{\topsep}{0pt}%
%                   \setlength{\parskip}{0pt}%
}%
}{%
\end{list}}


\newcommand{\ini}{\mathrm{s}}
\newcommand{\tar}{\mathrm{t}}

\newcommand{\X}{\mathcal{X}}
\newcommand{\Z}{\mathcal{Z}}
\newcommand{\sets}{\mathcal{S}}
\newcommand{\zind}{\mathcal{Z}_{\mathrm{ind}}}
\newcommand{\zsol}{\mathcal{Z}_{\mathrm{sol}}}
\newcommand{\labeln}{\mathsf{label}}
\newcommand{\rootn}{\mathsf{root}}
\newcommand{\child}[1]{\mathsf{child}_{#1}}
\newcommand{\remove}{\mathsf{remove}}
\newcommand{\add}{\mathsf{add}}
\newcommand{\swap}{\mathsf{swap}}

\newcommand{\OrigDecom}{\textsf{orig}\xspace}
\newcommand{\HeurDecom}{\textsf{heur}\xspace}
\newcommand{\OptDecom}{\textsf{opt}\xspace}

\newcommand{\TJ}{\textsc{Token Jumping}\xspace}
\newcommand{\satisfiability}{\textsc{SAT}\xspace}
\newcommand{\SR}{\textsc{SAT Reconfiguration}\xspace}
\newcommand{\LCR}{\textsc{List Coloring Reconfiguration}\xspace}
\newcommand{\SPR}{\textsc{Shortest Path Reconfiguration}\xspace}

\newcommand{\SATseries}{\textsc{SAT} series\xspace}
\newcommand{\LCseries}{\textsc{LC} series\xspace}
\newcommand{\SPseries}{\textsc{SP} series\xspace}
\newcommand{\ISseries}{\textsc{IS} series\xspace}

\newcommand{\XSAT}{\textsc{SAT}\xspace}
\newcommand{\XLC}{\textsc{LC}\xspace}
\newcommand{\XSP}{\textsc{SP}\xspace}
\newcommand{\XIS}{\textsc{IS}\xspace}
\newcommand{\XRAND}{\textsc{RAND}\xspace}

\newcommand{\snumfun}[1]{s_{#1}}
\newcommand{\snum}[2]{\snumfun{#1}(#2)}

\newcommand{\KJ}[2]{{\color{red} #2}}
\newcommand{\KJQ}[1]{{\color{blue} #1}}

\begin{document}

\title{Accelerating ZDD-Based Solvers for Independent Set Reconfiguration via Path Decomposition}
%\thanks{This work was supported by JSPS KAKENHI Grant Numbers JP24H00686, JP24H00690, and JP23K12974.}}

\author{
\IEEEauthorblockN{Kanae Sawamura}
\IEEEauthorblockA{\textit{Graduate School of Information Sciences}\\
\textit{Tohoku University}\\
Sendai, Japan\\
sawamura.kanae.r7@dc.tohoku.ac.jp}
\and
\IEEEauthorblockN{Takehide Soh}
\IEEEauthorblockA{\textit{Graduate School of Informatics}\\
\textit{Nagoya University}\\
Nagoya, Japan\\
soh@i.nagoya-u.ac.jp}
\and
\IEEEauthorblockN{Yasuaki Kobayashi}
\IEEEauthorblockA{\textit{Faculty of Information Science and Technology}\\
\textit{Hokkaido University}\\
Sapporo, Japan\\
koba@ist.hokudai.ac.jp}
\and
\IEEEauthorblockN{Yuma Tamura}
\IEEEauthorblockA{\textit{Graduate School of Information Sciences}\\
\textit{Tohoku University}\\
Sendai, Japan\\
tamura@tohoku.ac.jp}
\and
\IEEEauthorblockN{Yuta Nozaki}
\IEEEauthorblockA{\textit{Faculty of Science}\\
\textit{Hokkaido University}\\
Sapporo, Japan\\
\textit{WPI-SKCM$^2$, Hiroshima University}\\
Hiroshima, Japan\\
nozaki@math.sci.hokudai.ac.jp}\and
\IEEEauthorblockN{Jun Kawahara}
\IEEEauthorblockA{\textit{Graduate School of Informatics}\\
\textit{Kyoto University}\\
Kyoto, Japan\\
jkawahara@i.kyoto-u.ac.jp}
\and
\IEEEauthorblockN{Takehiro Ito}
\IEEEauthorblockA{\textit{Graduate School of Information Sciences}\\
\textit{Tohoku University}\\
Sendai, Japan\\
takehiro@tohoku.ac.jp}
}

\maketitle

\begin{abstract}
Zero-suppressed binary decision diagrams (ZDDs) are a standard data structure for
compactly representing families of combinatorial solutions, and reducing the
width of the path decomposition used to construct a ZDD is known to be effective
for keeping it small.
In combinatorial reconfiguration, however, a ZDD-based solver must build a
solution-space ZDD and then repeatedly perform transition computations on top of
it, and how much this width affects the end-to-end performance of such solvers has
remained unclear.
This paper studies path decomposition as a preprocessing step for the ZDD-based
independent set reconfiguration (ISR) solver \textsc{ddreconf} under the token
jumping rule.
By reordering the input edges according to a computed path decomposition, we
shrink the solution-space ZDD and improve the overall performance of
\textsc{ddreconf}, without modifying the solver itself.
Computational experiments on two variants of the problem, the shortest and the
farthest variant, show that the preprocessing consistently reduces ZDD size,
memory usage, and running time, and increases the number of instances solved
within the time limit.
Moreover, computing an optimal path decomposition is unnecessary in practice: a
fast heuristic decomposition attains essentially the same width and the same
gains, and on large instances, for which the optimal pathwidth cannot be computed
in reasonable time, the heuristic is the only option that lets the solver reach
them.
\end{abstract}

\begin{IEEEkeywords}
Combinatorial reconfiguration, independent set reconfiguration, zero-suppressed binary decision diagram, path decomposition, pathwidth
\end{IEEEkeywords}

\section{Introduction}
Combinatorial reconfiguration studies how one feasible solution of a
combinatorial problem can be transformed into another through a sequence of small
local changes~\cite{nishimura18}.
Among these problems, independent set reconfiguration (ISR) asks whether one
independent set of a graph can be transformed into another by a sequence of small
changes that keeps an independent set at every step.
In this paper we adopt the token jumping rule, in which each step moves a single
token to another vertex, equivalently, removes one vertex from the current
independent set and adds one vertex outside it.
ISR is PSPACE-complete in general~\cite{kamemi12,itdehapasiueun11}, and developing
practically efficient solvers remains an active research topic.

Zero-suppressed binary decision diagrams (ZDDs) provide a compact implicit
representation of a family of sets and have been successfully applied to a wide
range of combinatorial problems.
It is well known that the size of a ZDD built from a graph is closely related to
the pathwidth of the graph: a smaller pathwidth often leads to a smaller
ZDD~\cite{inoue16}.
This observation has motivated path-decomposition-based techniques in many ZDD
applications, where reordering the variables according to a good path
decomposition keeps the constructed ZDD small.

ZDD-based solvers for reconfiguration problems, however, follow a different
workflow.
They first construct a ZDD representing the entire solution space and then
repeatedly execute transition operations on this ZDD to search for a
reconfiguration sequence.
Although a path decomposition of small width is known to help the first step,
the construction of the solution-space ZDD, it has remained unclear how much it
helps the solver as a whole, once the cost of the subsequent transition
computation and the cost of computing the decomposition itself are taken into
account.

In this paper, we study path decomposition as a preprocessing step for the
state-of-the-art ZDD-based ISR solver \textsc{ddreconf}~\cite{ItoKNSSTT23} under
the token jumping rule.
The input to \textsc{ddreconf} is a graph given as an ordered list of edges, and
the solver processes these edges in the given order; this order determines the
width of the path decomposition that the solver implicitly works with, and hence
the sizes of the ZDDs it builds.
Our preprocessing simply reorders the edges according to a path decomposition of
small width, leaving the solver itself untouched.
We evaluate this idea on two variants that probe the solution space in
complementary ways: the shortest variant, which seeks a shortest reconfiguration
sequence between two given independent sets, and the farthest variant, which
seeks an independent set farthest from a given one.
Our main contributions are as follows:
\begin{itemize}
\item We show, through computational experiments, that reordering the input edges
by a path decomposition consistently improves \textsc{ddreconf}: it shrinks the
solution-space ZDD and reduces running time and memory, and it increases the
number of instances solved within the time limit, on both the shortest and the
farthest variant.
\item We show that computing an optimal path decomposition is unnecessary in
practice. A fast heuristic decomposition attains essentially the same width as
the optimal one and yields comparable gains; furthermore, on large instances
whose optimal pathwidth cannot be computed in reasonable time, the heuristic is
the only one of the two that lets the solver reach them.
\item We find that the width of the decomposition used by the solver is a strong
predictor of whether an instance can be solved, more so than the size of the
graph.
\end{itemize}

The remainder of this paper is organized as follows.
Section~\ref{sec:preliminaries} introduces ISR and reviews ZDD-based solvers.
Section~\ref{sec:approach} presents our path-decomposition-based preprocessing
and its integration with \textsc{ddreconf}.
Section~\ref{sec:experiments} reports experimental results.
Section~\ref{sec:conclusion} concludes the paper.


\section{Preliminaries}
\label{sec:preliminaries}

\subsection{Independent Set Reconfiguration}
\label{subsec:isr}

Let $G=(V,E)$ be an undirected graph, where $V$ is the set of vertices and $E$
is the set of edges. A vertex subset $I\subseteq V$ is an \emph{independent set
of $G$} if no two vertices in $I$ are adjacent in $G$.
Henceforth, we refer to each vertex contained in an independent set as a
\emph{token}.
Two independent sets $I$ and $I'$ of $G$ are \emph{adjacent} under the
\emph{Token Jumping} (TJ) rule if $|I \bigtriangleup I'| = 2$, namely, $I'$ can
be obtained from $I$ by moving a single token in $I$ to another vertex that is
not in $I$, and vice versa.
Then, for two independent sets $I_\ini$ and $I_\tar$ of $G$ such that
$|I_\ini|=|I_\tar|$, a \emph{reconfiguration sequence} between $I_\ini$ and
$I_\tar$ is a sequence $\langle I_0, I_1, \ldots, I_\ell \rangle$ such that:
\begin{listing}{-}
\item[-] $I_0 = I_\ini$ and $I_\ell = I_\tar$;
\item[-] $I_i$ is an independent set of $G$ such that $|I_i| = |I_\ini| = |I_\tar|$, for each $i$, $0 \leq i \leq \ell$;
\item[-] $I_i$ and $I_{i+1}$ are adjacent under TJ for each $i$, $0 \le i \le \ell-1$.
\end{listing}
Then, the \emph{Independent Set Reconfiguration Problem} (ISRP) under TJ is to
determine whether there exists a reconfiguration sequence between $I_\ini$ and
$I_\tar$. This problem is known to be \PSPACE-complete~\cite{kamemi12}.
The \PSPACE-completeness implies that there are yes-instances such that even a
shortest reconfiguration sequence has super-polynomial length, under the
assumption that $\PSPACE \neq \NP$.

In this paper, we study two variants of ISRP under TJ, which probe the structure
of the solution space in complementary ways.
The \emph{solution space} of an ISR instance $(G, I_\ini, I_\tar)$ is the graph whose nodes are the
independent sets of $G$ of size $|I_\ini|$ and in which two independent sets are
joined whenever they are adjacent under TJ; a reconfiguration sequence is exactly
a path in this graph. The \emph{distance} between two independent sets in the
solution space is the length of a shortest reconfiguration sequence between them;
or it is $\infty$ if no such sequence exists. The solution space can be exponentially
large in $|V|$, which makes designing practically efficient solvers a challenging
task.

The first is the \emph{shortest variant}, which generalizes ISRP.
Given an ISRP instance $(G, I_\ini, I_\tar)$, the shortest variant aims to
determine the existence of a reconfiguration sequence between $I_\ini$ and
$I_\tar$ and, if so, to compute the distance between them. This is the classical
reachability/shortest-path question between two prescribed independent sets in the solution space,
which has been studied extensively from a theoretical standpoint and was adopted
as a metric of the CoRe Challenge~\cite{sotaokit24}.

The second is the \emph{farthest variant}~\cite{ItoKNSSTT23}.
Given only one independent set $I_\ini$, the farthest variant aims to find an
independent set whose distance from $I_\ini$ is maximum among all independent
sets reachable from $I_\ini$, together with that distance. Rather than connecting
two prescribed independent sets, this variant asks how far the reconfiguration
relation can carry $I_\ini$ throughout its connected component of the solution
space. It is therefore a one-to-many question about the extent of a connected
component, in contrast to the one-to-one, path-between-two-endpoints nature of
the shortest variant. Studying both lets us observe how the width of the decomposition affects a
ZDD-based solver on two qualitatively different uses of the same solution-space
representation.

\subsection{Overview of ZDD-Based Solvers}
\label{subsec:zdd_solver}

A \emph{zero-suppressed binary decision diagram} (ZDD) is a directed acyclic
graph that compactly represents a family of sets over a fixed universal set
$U = \{x_1, \dots, x_n\}$ equipped with a total order $x_1 < \cdots < x_n$.
It has two terminal nodes $\bot$ and $\top$, representing the empty family
$\emptyset$ and the family $\{\emptyset\}$, respectively, and each non-terminal
node is labeled with an element of $U$ and has a $0$-arc and a $1$-arc.
By sharing isomorphic subgraphs and suppressing redundant nodes, a ZDD can
represent a family whose cardinality is exponential in $n$ using far fewer
nodes, and it supports efficient family operations such as union,
intersection, and difference through recursive procedures over its shared
structure~\cite{Minato1993,knuth:taocp41}.

For ISRP, the relevant family is the set of all independent sets of the input
graph $G = (V, E)$, where the universal set is taken to be the vertex set $V$.
The state-of-the-art ZDD-based solver \textsc{ddreconf}~\cite{ItoKNSSTT23} solves
ISRP under the TJ rule in two phases.

In the first phase, the solver constructs the \emph{solution-space ZDD}.
It first builds a ZDD $\zind$ representing the family of all independent sets of
$G$~\cite{Hayase1995}. Recall that, under the TJ rule, every independent set in a
reconfiguration sequence has the same cardinality $k = |I_\ini|$ as the start
independent set $I_\ini$. The solver therefore builds a ZDD $\Z_{=k}$
representing all vertex subsets of cardinality exactly $k$~\cite{knuth:taocp41},
and obtains the solution space as $\zsol = \zind \cap \Z_{=k}$, which represents
every independent set of $G$ of cardinality $k$. The family $\sets(\zsol)$ is
exactly the node set of the solution space defined in
Section~\ref{subsec:isr}.

In the second phase, the solver repeatedly performs \emph{transition
computation} on top of $\zsol$. Starting from $\Z^{0}$ with
$\sets(\Z^{0}) = \{I_\ini\}$, it iteratively computes
\begin{equation}
  \Z^{i} \;\gets\; \mathsf{op}(\Z^{i-1}) \cap \zsol,
  \label{eq:transition}
\end{equation}
where $\mathsf{op}$ applies the TJ rule to every independent set encoded by
$\Z^{i-1}$. Concretely, $\mathsf{op}(\Z) = \swap(\Z, V, V)$, where
$\swap(\Z, A, R)$ represents the family
\[
  \{\, I \cup \{v\} \setminus \{v'\} \mid I \in \sets(\Z),\,
  v \in A \setminus I,\, v' \in I \cap R \,\}
\]
obtained by moving a single token, and intersecting with $\zsol$ retains only
those sets that are independent sets of cardinality $k$. Hence $\Z^{i}$
represents exactly the family of independent sets whose distance from $I_\ini$ is
$i$, and iterating \eqref{eq:transition} amounts to a breadth-first search over
the solution space. This implicit search lets \textsc{ddreconf} explore an
exponentially large solution space, since each $\Z^{i}$ may encode an enormous
number of independent sets within a small ZDD.

Both variants are answered from this sequence of ZDDs. For the shortest variant,
the solver checks $\Z^{i}$ after each step: once $I_\tar \in \sets(\Z^{i})$, the
distance between $I_\ini$ and $I_\tar$ is $i$, and once $\Z^{i} = \bot$ without
having reached $I_\tar$, no reconfiguration sequence exists. For the farthest
variant, the solver continues until $\Z^{h+1} = \bot$; the last nonempty
$\Z^{h}$ then contains the independent sets farthest from $I_\ini$, and $h$ is
their distance.

The performance of this workflow is governed by the \emph{size} (number of
nodes) of the ZDDs that appear in it, which in turn depends critically on the
order in which the edges of $G$ are processed during ZDD construction. For the
first phase, this dependence is well understood~\cite{inoue16}: the working state
maintained during top-down ZDD construction, and hence both the construction
time and the resulting ZDD size, grows rapidly with the width of the path
decomposition induced by the processing order (we make this correspondence
precise in Section~\ref{subsec:pathdecomp}). A processing order of small width
therefore keeps $\zind$, and hence $\zsol$, compact; for example, the family of
independent sets of an $8 \times 250$ grid graph has cardinality
$3.07 \times 10^{361}$, yet it can be represented by a ZDD with only $49{,}989$
nodes under a suitable order~\cite{ItoKNSSTT23}. Crucially, \textsc{ddreconf}
fixes a single processing order before construction and reuses it throughout
both phases, so this order influences not only the construction of $\zsol$ but
also the entire subsequent transition computation. While the benefit of a
small-width order for the first phase is inherited from established results on
ZDD construction, whether it improves the solver as a whole, once the cost of
the transition computation and the cost of obtaining the order itself are taken
into account, is the question we study in this paper.

\section{Our Approach}
\label{sec:approach}

Our approach is built on a simple observation made in
Section~\ref{subsec:zdd_solver}: the performance of \textsc{ddreconf} is governed
by the sizes of the ZDDs that appear throughout its two phases, and these sizes
are determined by the order in which the edges of the input graph $G$ are
processed during ZDD construction. We already know that, for the first phase, a
processing order of small width keeps the constructed ZDD small; what is not
clear is whether this benefit carries over to the solver as a whole, once the
cost of the subsequent transition computation and the cost of obtaining the
order itself are taken into account. The original \textsc{ddreconf} fixes the
order according to the input itself and never revisits it. We instead treat the
order as a tunable \emph{preprocessing} parameter derived from a path
decomposition of $G$, and feed the resulting order to \textsc{ddreconf} without
modifying the solver. Concretely, we compare three decompositions of increasing quality
and preprocessing cost: \OrigDecom, the decomposition induced by the input edge
order; \HeurDecom, a heuristic decomposition; and \OptDecom, a minimum-width
decomposition. This
section formalizes the path-decomposition view of the input order
(Section~\ref{subsec:pathdecomp}), describes the heuristic and exact algorithms
behind \HeurDecom and \OptDecom (Section~\ref{subsec:algorithms}), and explains
how the preprocessing is integrated into the \textsc{ddreconf} pipeline
(Section~\ref{subsec:integration}).

\subsection{Path Decomposition and the Input Edge Order}
\label{subsec:pathdecomp}

We first recall the standard definitions.
A \emph{path decomposition} of a graph $G = (V, E)$ is a sequence
$(X_1, \dots, X_\ell)$ of subsets of $V$, called \emph{bags}, such that
(i) for every edge $\{u, v\} \in E$ there is a bag $X_i$ with $u, v \in X_i$, and
(ii) for every vertex $v \in V$, the bags containing $v$ form a contiguous
subsequence. The \emph{width} of a path decomposition is
$\max_{1 \le i \le \ell} |X_i| - 1$, and the \emph{pathwidth} $\mathrm{pw}(G)$ of
$G$ is the minimum width over all its path decompositions.

The CoRe Challenge provides each instance as an input graph in the DIMACS
format~\cite{sotaokit24}: a \texttt{p} line declares the numbers of vertices and
edges, and each subsequent \texttt{e} line lists the two endpoints of an edge.
When \textsc{ddreconf} reads such a file and constructs $\zind$ in the first
phase, it processes the edges one by one, in the order in which they appear in
the file. The cost of this top-down construction is governed by the
\emph{frontier}, the set of vertices that are incident to both already-processed
and not-yet-processed edges; it is a folklore fact, made explicit by Inoue and
Minato~\cite{inoue16}, that the maximum frontier size over the construction
coincides with the width of a path decomposition of $G$ determined by the edge
order. Moreover, given an edge order, the width of this associated path
decomposition can be computed by the algorithm of Kobayashi and
Nakahata~\cite[Lemma~2]{KobayashiNakahata2020}. Through this correspondence, each edge order has a well-defined \emph{width},
the width of its associated path decomposition; this width is at least
$\mathrm{pw}(G)$, with equality exactly when the order is optimal.

Consequently, the order in which edges are listed in the input file, an artifact
of how each instance happens to be encoded, directly determines the width of the
path decomposition that \textsc{ddreconf} implicitly works with, and hence the
sizes of the ZDDs it builds and manipulates. For the first phase, the effect is
well understood: a smaller width keeps $\zind$, and hence $\zsol$, compact
(Section~\ref{subsec:zdd_solver}). Whether this carries over to the end-to-end
performance of the solver is the subject of this paper. This leads to a natural
plan: \emph{replace the input edge order with one derived from a path
decomposition of small width, and measure how the overall performance of
\textsc{ddreconf} changes.}

\subsection{Heuristic and Optimal Path Decomposition}
\label{subsec:algorithms}

Computing $\mathrm{pw}(G)$ exactly is NP-hard, so obtaining a minimum-width path
decomposition is expensive in general. Besides the decomposition induced by the
input order (\OrigDecom), we obtain a path decomposition using the algorithm of
Kobayashi, Komuro, and Tamaki~\cite{KobayashiKomuroTamaki2014}. This algorithm
computes an optimal path decomposition by reducing the search space through
commitments and gradually narrowing down to minimum width; it produces a first
feasible decomposition quickly and, if run to completion, eventually returns an
optimal one. We exploit both ends of this process, yielding two path
decompositions of different quality and cost. This algorithm is an existing
method, which we use as is and do not claim as a contribution of this paper.

\paragraph{Heuristic decomposition (\HeurDecom).}
As a fast preprocessing step, we take the first decomposition that the algorithm
produces, whose width is small but not necessarily minimum. Obtaining it is
quick enough to serve as a lightweight preprocessor even on the larger CoRe
Challenge instances. Because this first decomposition is not guaranteed to attain
the minimum width, it is occasionally wider than the one induced by the original
input order. Since the width of \OrigDecom is itself available
(Section~\ref{subsec:pathdecomp}), we compare the two and keep the better one:
only when the first decomposition is strictly wider than the original do we fall
back to the original order. Thus \HeurDecom never yields a wider decomposition
than \OrigDecom.

\paragraph{Optimal decomposition (\OptDecom).}
To examine the effect of attaining the true minimum width, we run the same
algorithm to completion, obtaining an optimal path decomposition, i.e., one of
width exactly $\mathrm{pw}(G)$. This is far more time-consuming than obtaining the
first decomposition, and is prohibitively so for some instances, but it lets us
measure the best possible gain achievable through path-decomposition-based
reordering and compare it against \HeurDecom.

\subsection{Integration with \textsc{ddreconf} and the Three Decompositions}
\label{subsec:integration}

In all three cases, the path decomposition is used purely as preprocessing: from
the decomposition we derive an edge order, write it back as a reordered DIMACS
file, and run the unmodified \textsc{ddreconf} on it. The solver itself is
identical across all runs, including its solution-space ZDD construction, its
swap-based transition computation, and its breadth-first reconfiguration search;
only the order in which the edges of $G$ are presented differs. This makes the
comparison clean: any difference in running time, memory usage, or
solved-instance count is attributable solely to the edge order, and hence to the
underlying path decomposition.

We compare the following three path decompositions.
\begin{description}
  \item[\OrigDecom:] the path decomposition induced by the edge order as given in
    the CoRe Challenge input file, used without any reordering. This is the
    baseline and corresponds to the original \textsc{ddreconf} behavior.
  \item[\HeurDecom:] a path decomposition of small but not necessarily minimum
    width, computed by the heuristic algorithm~\cite{KobayashiKomuroTamaki2014};
    if it is strictly wider than \OrigDecom, we fall back to the original order,
    so that \HeurDecom is never wider than \OrigDecom.
  \item[\OptDecom:] an optimal (minimum-width) path decomposition, computed by an
    exact algorithm.
\end{description}
For \HeurDecom and \OptDecom, we derive an edge order from the computed path
decomposition and run \textsc{ddreconf} on that order. The preprocessing cost
increases sharply along this list: \OrigDecom incurs none, \HeurDecom adds only
a fast preprocessing step, whereas \OptDecom can be extremely time-consuming.
Section~\ref{sec:experiments} reports how these three path decompositions compare
in practice, and in particular whether the inexpensive \HeurDecom decomposition
already captures most of the benefit that the costly \OptDecom decomposition
provides.


\section{Experimental Results}
\label{sec:experiments}

\subsection{Experimental Setup}
\label{subsec:setup}

We implemented the path-decomposition preprocessing on top of
\textsc{ddreconf}~\cite{ItoKNSSTT23} and evaluated the three decompositions
\OrigDecom, \HeurDecom, and \OptDecom defined in Section~\ref{sec:approach}.
%% TODO: machine spec (CPU, RAM), compiler, ZDD libraries and versions,
%% ddreconf version/commit, and the implementation used for the pathwidth
%% computation of Kobayashi--Komuro--Tamaki.
All experiments were run on \todo{machine and software environment}.

\paragraph{Benchmarks.}
We used the benchmark instances of the CoRe Challenge 2022 and
2023~\cite{sotaokit24}. In total there are $693$ instances arising from $519$
distinct input graphs; each graph may appear in several instances that differ in
the start and goal independent sets.
%% TODO: how instances were selected (all? a subset?), and how many of the 519
%% graphs are too large to be read by ddreconf at all (handoff item 9). Under
%% \HeurDecom, ddreconf rejects 73 instances outright; clarify this in the
%% population description.
The graphs range widely in size, from tens of vertices to graphs with more than
$10^{6}$ edges \todo{confirm size range and selection criteria}.

\paragraph{Time limits.}
We adopt a two-stage time budget that separates preprocessing from solving.
Computing \OptDecom requires solving an NP-hard pathwidth problem, so we allow up
to one hour for this computation; if it does not finish within the limit, no
optimal decomposition is available for that instance and \OptDecom is undefined.
Computing \HeurDecom and reading the input order for \OrigDecom incur negligible
preprocessing cost.
Given the (possibly reordered) input, we then run the unmodified
\textsc{ddreconf} with a time limit of $30$ minutes of wall-clock time per
instance. An instance is counted as \emph{solved} under a given decomposition if
\textsc{ddreconf} terminates with the correct answer within this $30$-minute
limit.
We report results for both the shortest and the farthest variant.

Of the $519$ graphs, the optimal pathwidth (and hence \OptDecom) was computed
within the one-hour limit for $226$ graphs, corresponding to $279$ of the $693$
instances; for the remaining $293$ graphs ($414$ instances) the optimal decomposition
could not be obtained.
This split is itself informative: the instances for which \OptDecom is available
are markedly smaller (median $120$ vertices and $476$ edges) than those for which
it is not (median $450$ vertices and $1560$ edges). We therefore present two
kinds of comparison: one restricted to the instances on which all three decompositions
are available, and one over all instances that contrasts \OrigDecom with
\HeurDecom.

\subsection{Width of the Three Decompositions}
\label{subsec:pw-three}

Before turning to solver performance, we examine how much the three decompositions differ
in width on the input graphs themselves, independently of whether the instances
are solved.
Figure~\ref{fig:pw-three} shows, for each of the $226$ graphs whose optimal
pathwidth was computed within the one-hour limit, the width of \OrigDecom,
\HeurDecom, and \OptDecom, with the three values for the same graph connected by
a vertical line.
Two observations stand out.
First, reordering the edges by a path decomposition reduces the width
dramatically: the median width drops from $32$ for \OrigDecom to $11$ for
\HeurDecom and $10$ for \OptDecom (the last value is the pathwidth
$\mathrm{pw}(G)$).
Second, \HeurDecom is already very close to \OptDecom; the two markers nearly
coincide for most graphs, and the median per-graph ratio of their widths is
$1.00$.
In other words, on this benchmark the heuristic decomposition attains
essentially the same width as the optimal one, while costing far less to compute
(Section~\ref{subsec:solved}).
Many graphs whose width is on the order of $100$ under the input order collapse
to a narrow band of width around $10$ to $17$ once decomposed, even when the
graph is large.

\begin{figure}[t]
  \centering
  %% TODO: final vector figure; recolor to avoid clashing with the
  %% solved/unsolved colors used in Fig.~\ref{fig:pw-solvability}.
  \includegraphics[width=\linewidth]{pw_three_orders.pdf}
  \caption{Width of \OrigDecom, \HeurDecom, and \OptDecom for each input
  graph (the $226$ graphs whose optimal pathwidth was computed within one hour).
  The three values for the same graph are joined by a vertical line.
  \HeurDecom (square) and \OptDecom (triangle) almost coincide, well below
  \OrigDecom (circle).}
  \label{fig:pw-three}
\end{figure}


\subsection{Solved Instances and Running Time}
\label{subsec:solved}

We now turn to the bottom-line question of how many instances each decomposition lets
\textsc{ddreconf} solve within the $30$-minute limit.
Table~\ref{tab:solved} reports the counts for both variants, under the two
comparisons motivated in Section~\ref{subsec:setup}.

\begin{table}[t]
  \centering
  \caption{Number of instances solved within $30$ minutes. The left block is
  restricted to the $279$ instances for which all three decompositions are available; the
  right block is over all $693$ instances, where \OptDecom is available for only
  $279$ of them.}
  \label{tab:solved}
  \begin{tabular}{lrrr|rr}
    \hline
    & \multicolumn{3}{c|}{opt-available ($279$)} & \multicolumn{2}{c}{all ($693$)} \\
    variant & \OrigDecom & \HeurDecom & \OptDecom & \OrigDecom & \HeurDecom \\
    \hline
    shortest & $210$ & $216$ & $218$ & $210$ & $231$ \\
    farthest & $208$ & $215$ & $218$ & $208$ & $229$ \\
    \hline
  \end{tabular}
\end{table}

On the instances where all three decompositions are available, both \HeurDecom and
\OptDecom solve more than \OrigDecom, and \HeurDecom is nearly on par with
\OptDecom (for the shortest variant, $216$ versus $218$; for the farthest
variant, $215$ versus $218$).
Over all instances, \HeurDecom solves the most: $231$ versus $210$ for
\OrigDecom on the shortest variant, and $229$ versus $208$ on the farthest
variant.

The reason \HeurDecom is preferable to \OptDecom in practice lies in
preprocessing cost.
On the instances where \OptDecom is available, computing \HeurDecom takes a
median of $2.1$ seconds (at most $13.7$ seconds), whereas computing \OptDecom
takes a median of $3.2$ seconds but a mean of $64.7$ seconds and up to
$2348$ seconds.
More importantly, for $414$ of the $693$ instances the optimal pathwidth could
not be computed within one hour at all, so \OptDecom never reaches the solver on
these instances.
These are exactly the large, hard instances (median $450$ vertices and $1560$
edges).
On these instances, \HeurDecom lets \textsc{ddreconf} solve $15$ instances on the
shortest variant and $14$ on the farthest variant that it would otherwise not
reach.
To confirm that this is a genuine effect of the heuristic decomposition rather than of
the instances merely being easy, we ran \OrigDecom on these heuristic-solved
instances as well.
%% TODO: report the final count once the in-progress \OrigDecom runs on the
%% 15 (shortest) / 14 (farthest) heur-solved large instances complete. Several
%% are already confirmed to exceed the 30-minute limit under \OrigDecom while
%% \HeurDecom solves them in seconds (e.g.\ width drops $234\to23$ and
%% $140\to17$); a few have unchanged width under reordering and must be
%% reported separately or excluded (handoff: groups A/D vs C).
\OrigDecom fails to solve \todo{N} of them within the $30$-minute limit, whereas
\HeurDecom solves them in seconds.
Thus, on the large instances where an optimal decomposition is out of reach, the
lightweight heuristic decomposition is the one that lets the solver make
progress.

Because the solution-space ZDD construction is shared by both variants, the two
variants differ only in the second phase, where the farthest variant explores the
solution space until it is exhausted.
This makes the farthest variant somewhat more expensive in the second phase, and
correspondingly a few fewer instances are solved under it; otherwise the two
variants exhibit the same trends.

\subsection{Decomposition Width Determines Solvability}
\label{subsec:solvability}

Finally, we examine which property of an instance governs whether it can be
solved.
Figure~\ref{fig:pw-solvability} plots every instance by its number of edges and
the width of the decomposition used, marking whether \textsc{ddreconf} solves it
within $30$ minutes, for \OrigDecom (top) and \HeurDecom (bottom).
In both panels, solved and unsolved instances separate along an almost horizontal
boundary: solvability is governed by this width rather than by graph size.
Large graphs are solved when this width is small and left unsolved when it
is large; among instances with more than the median number of edges, those solved
under \HeurDecom have a median width of $17$, while those left unsolved have a
median width of $156$.
Comparing the two panels shows the effect of preprocessing: \HeurDecom pulls the
same graphs downward into the low-width, solvable region, turning many
instances that are unsolved under \OrigDecom into solved ones.
Quantitatively, on a logarithmic scale matching the axes of the figure, the
point-biserial correlation between solvability and width is stronger than
that between solvability and the number of edges ($-0.64$ versus $-0.53$ for
\HeurDecom), confirming that width is the better predictor.

\begin{figure}[t]
  \centering
  %% TODO: final vector figure; the farthest variant shows the same trend and is
  %% omitted to save space (state this if a reviewer asks).
  \includegraphics[width=\linewidth]{pw_solvability.pdf}
  \caption{Each instance plotted by number of edges and the width of its
  decomposition, for \OrigDecom (top) and \HeurDecom (bottom), on the shortest variant.
  Solved instances (within $30$ minutes) lie below an almost horizontal boundary,
  showing that width, not size, governs solvability. \HeurDecom shifts
  instances downward, enlarging the solvable region.}
  \label{fig:pw-solvability}
\end{figure}



\section{Analyses}

\subsection{Effect on the Solution-Space ZDD}
\label{subsec:zddsize}

We next measure how the reduced width translates into a smaller solution-space
ZDD, i.e., the ZDD $\zsol$ built in the first phase of \textsc{ddreconf}.
To compare the three decompositions on a common footing, we restrict attention to the
instances solved by all three decompositions within the time limit ($210$ instances for
the shortest variant and $208$ for the farthest variant), so that every decomposition
produces a value for every instance.
%% TODO: revisit whether to restrict to commonly solved instances (handoff
%% item 4-4). The maximum is now reported alongside the median.
Table~\ref{tab:zddsize} reports, for each decomposition and each variant, the
median and the maximum of the width, the solution-space ZDD size, the first-phase
construction time, and the peak memory, over those instances.

\begin{table*}[t]
  \centering
  \caption{Width and cost on the instances that all three decompositions solve
  within the time limit ($210$ for the shortest variant, $208$ for the farthest).
  For each quantity we report the geometric mean and the maximum over those
  instances: the decomposition width, the solution-space ZDD size (number of
  nodes), the first-phase ZDD construction time, the total wall-clock solving
  time, and the peak resident memory.}
  \label{tab:zddsize}
  \setlength{\tabcolsep}{4pt}
  \begin{tabular}{ll|rr|rr|rr|rr|rr}
    \hline
    & & \multicolumn{2}{c|}{width} & \multicolumn{2}{c|}{ZDD size}
      & \multicolumn{2}{c|}{constr.\ time [s]} & \multicolumn{2}{c|}{total time [s]}
      & \multicolumn{2}{c}{memory [MB]} \\
    variant & decomp. & geo. & max & geo. & max & geo. & max
      & geo. & max & geo. & max \\
    \hline
    \multirow{3}{*}{shortest}
      & \OrigDecom & $34.5$ & $611$ & $3{,}011$ & $142{,}539{,}591$ & $0.183$ & $1497.2$ & $3.04$ & $1711.7$ & $282$ & $137{,}248$ \\
      & \HeurDecom & $12.5$ & $485$ & $1{,}191$ & $793{,}894$ & $0.078$ & $1482.4$ & $1.44$ & $1715.0$ & $170$ & $134{,}664$ \\
      & \OptDecom  & $11.7$ & $485$ & $971$ & $395{,}153$ & $0.072$ & $1459.2$ & $1.28$ & $1706.8$ & $158$ & $134{,}662$ \\
    \hline
    \multirow{3}{*}{farthest}
      & \OrigDecom & $34.3$ & $611$ & $2{,}805$ & $142{,}539{,}591$ & $0.159$ & $1475.9$ & $2.94$ & $1629.5$ & $300$ & $137{,}247$ \\
      & \HeurDecom & $12.4$ & $485$ & $1{,}157$ & $793{,}894$ & $0.074$ & $1231.5$ & $1.66$ & $1687.6$ & $184$ & $134{,}664$ \\
      & \OptDecom  & $11.6$ & $485$ & $948$ & $395{,}153$ & $0.070$ & $1229.3$ & $1.46$ & $1683.7$ & $165$ & $134{,}663$ \\
    \hline
  \end{tabular}
\end{table*}

While a typical instance yields a small solution-space ZDD under every
decomposition, so that the medians are comparable, the benefit of a narrow
decomposition is decisive on the hardest instances: the largest solution-space
ZDD falls from over $1.4\times10^{8}$ nodes under \OrigDecom to about
$7.9\times10^{5}$ under \HeurDecom and $4.0\times10^{5}$ under \OptDecom, with a
comparable gap in peak memory.
For the shortest variant, relative to \OrigDecom the per-instance reduction in
ZDD size has a geometric mean of $2.53\times$ for \HeurDecom and $3.10\times$ for
\OptDecom; the corresponding reductions in first-phase construction time and peak
memory are $2.4\times$ and $1.7\times$ for \HeurDecom, and $2.5\times$ and
$1.8\times$ for \OptDecom. The farthest variant shows the same pattern
($2.42\times$ and $2.96\times$ in ZDD size for \HeurDecom and \OptDecom).
Across these commonly solved instances, smaller width is consistently associated
with a smaller solution-space ZDD: the Spearman rank correlation between
width and ZDD size is $0.66$, and that between ZDD size and total solving
time is $0.51$.
As in the width comparison, \HeurDecom captures most of the benefit of
\OptDecom: it attains essentially the same width (median per-graph ratio $1.00$)
and about $80\%$ of the optimal geometric-mean reduction in the size of the
solution-space ZDD.

%% =====================================================================
%% Draft subsection: Path Decomposition in the Second Phase
%% Figures referenced: fig_secondphase_trajectory.pdf, fig_secondphase_scatter.pdf
%% Numbers computed from the 29-instance per-step CSV (node counts, hence
%% independent of the time limit; the 30 vs 40 min difference does not affect
%% any statement below).
%% =====================================================================
\subsection{Path Decomposition in the Second Phase}
\label{subsec:secondphase}

The width of a decomposition is known to control the size of the solution-space
ZDD $\zsol$ built in the first phase. It is far less obvious whether it also
controls the \emph{second} phase, in which a new ZDD $Z^i$ is constructed at
every reconfiguration step. To examine this we recorded, for a diverse set of
$29$ instances, the node count of every intermediate ZDD $Z^i$ under each of the
three decompositions. Because these node counts are determined by the
decomposition alone, the comparison below is independent of the time limit.

\paragraph{The first-phase benefit extends to the second phase.}
Whenever a narrower decomposition shrinks $\zsol$, it shrinks each intermediate
$Z^i$ by a comparable factor. Figure~\ref{fig:sp-traj}(left) shows a
representative case: on \textsf{mug88\_1} (widths $31/11/5$ for
\OrigDecom/\HeurDecom/\OptDecom) the peak intermediate ZDD reaches
$1.4\times10^{6}$ nodes under \OrigDecom but only $3.7\times10^{4}$ under
\HeurDecom and $5.0\times10^{3}$ under \OptDecom, a $281\times$ reduction, even
though all three runs complete. Comparable reductions of the peak intermediate
ZDD recur across the benchmark, for example $216\times$ on \textsf{miles500\_02},
$57\times$ on \textsf{huck}, and $15\times$ on \textsf{ph-04-03}. Over the
instances on which all three decompositions complete, the per-instance
second-phase reduction is strongly rank-correlated with the first-phase
reduction (Spearman $0.91$). The decomposition therefore acts on the whole
computation, not only on the construction of $\zsol$.

\paragraph{The reduction is sub-proportional.}
The second-phase reduction is consistently smaller than the first-phase
reduction (Figure~\ref{fig:sp-scatter}): on \textsf{miles500\_02} the
solution-space ZDD shrinks by a factor of $5.8\times10^{3}$ but the peak
intermediate ZDD only by $216$. This is expected, since each $Z^i$ represents
the wavefront of solutions reachable in exactly $i$ steps, a subset of $\zsol$
whose size is bounded independently of the decomposition. Even so, the
decomposition yields one to two orders of magnitude of second-phase savings on
the harder instances.

\paragraph{Both failure modes are relieved.}
Among the instances that \OrigDecom fails to solve within the time limit but
\HeurDecom or \OptDecom solves, the bottleneck lies in different phases
depending on the width. For high widths (\textsf{queen8\_12}, width $85$;
\textsf{games120}, width $91$) \OrigDecom cannot even build $\zsol$, so no
intermediate ZDD is ever produced: a first-phase failure. For moderate widths
(\textsf{mug}, widths $26$--$32$) \OrigDecom builds $\zsol$ quickly but the
intermediate ZDDs explode, reaching $1$--$2\times10^{6}$ nodes within the first
$14$--$15$ of about $20$ steps before timing out: a second-phase failure, whose
mechanism is confirmed by the completed sibling \textsf{mug88\_1} above. A narrow
decomposition removes the bottleneck wherever it occurs.

\paragraph{When width is uninformative.}
For a class of structured instances the pathwidth is a loose upper bound and
predicts neither phase. On \textsf{suzuki} (widths up to $611$) and
\textsf{DSJR500.1c} (width $496$) both $\zsol$ and every $Z^i$ remain small under
all three decompositions, and reordering by width yields no benefit
(Figure~\ref{fig:sp-traj}, right). An extreme case is \textsf{SAT\_exp\_012},
where $\zsol$ shrinks by $3.4\times10^{3}$ yet the peak intermediate ZDD is
essentially unchanged at about $1.1\times10^{3}$ nodes under every decomposition.
Across all runs and decompositions the correlation between the nominal width and
the peak intermediate ZDD size is consequently weak (Spearman $0.34$), even
though the within-instance effect is strong wherever the width genuinely binds.
This is why a second-phase trend is invisible in aggregate but clear once such
loose-bound instances are separated.

\paragraph{Long sequences.}
Finally, the second phase makes very long reconfiguration sequences tractable
precisely because the per-step ZDD is width-bounded. On the \textsf{hc-power}
family (width $5$) the intermediate ZDD stays near $5{,}000$ nodes across up to
$221{,}003$ steps, with the three decompositions coinciding; here the
reconfiguration length, not the per-step ZDD size, governs the total cost.

%% --------------------------------------------------------------------
\begin{figure}[t]
  \centering
  \includegraphics[width=\linewidth]{fig_secondphase_trajectory.pdf}
  \caption{Intermediate ZDD size $|Z^i|$ at each reconfiguration step $i$ (log
  scale, shortest variant). Left: \textsf{mug88\_1}, where the decompositions
  differ in width ($31/11/5$); the per-step ZDD is up to two orders of magnitude
  larger under \OrigDecom. Right: \textsf{suzuki\_2\_20}, whose pathwidth is a
  loose upper bound ($611$ vs $14$); the three curves coincide, so reordering by
  width gives no benefit.}
  \label{fig:sp-traj}
\end{figure}

\begin{figure}[t]
  \centering
  \includegraphics[width=0.82\linewidth]{fig_secondphase_scatter.pdf}
  \caption{Second-phase reduction (ratio of peak intermediate ZDD sizes,
  \OrigDecom over \OptDecom) against first-phase reduction (ratio of
  solution-space ZDD sizes), one point per instance, log--log. Points lie below
  the diagonal: the second-phase reduction tracks the first-phase reduction but
  is sub-proportional. The cluster near $(1,1)$ are instances whose pathwidth is
  a loose upper bound (\textsf{suzuki}, \textsf{DSJR}, \textsf{sp},
  \textsf{hc-power}); \textsf{SAT\_exp} is an extreme decoupling case.}
  \label{fig:sp-scatter}
\end{figure}

%\subsection{Reconfiguration Length}
%\label{subsec:reconflength}

A distinctive strength of the ZDD-based approach is its handling of long
reconfiguration sequences.
By \PSPACE-completeness, a shortest reconfiguration sequence may be
super-polynomially long (Section~\ref{subsec:isr}).
Because \textsc{ddreconf} advances its breadth-first search by one transition
layer at a time on top of a compressed solution-space ZDD, rather than unrolling
the sequence length explicitly, it can in principle reach such long sequences,
in contrast to bounded-model-checking solvers that increase the target length
one step at a time.
We examine this on our data.

Table~\ref{tab:reconflen} groups the instances by the length $\ell$ of the
reconfiguration sequence and reports, within the opt-available population and for
instances solved within $30$ minutes, how many each decomposition solves.
%% TODO: the population here is opt-available instances solved within 30 min, so
%% that the length is established under the same protocol; revisit if the setup
%% population description changes.
The instances span more than five orders of magnitude in $\ell$, from $1$ up to
about $2.2\times10^{5}$, and all are solved within the limit.
In the long-length ranges ($\ell>10^{4}$) the three decompositions solve exactly the
same instances, whereas the differences between them are concentrated in the
short-to-medium range, which corresponds to the wider graphs.
Within this population, \OptDecom solves every instance that any decomposition solves.
(On the shortest variant there are in addition $10$ instances with no
reconfiguration sequence, all of which are resolved by all three decompositions; the
farthest variant has none.)

\begin{table}[t]
  \centering
  \caption{Number of instances solved within $30$ minutes, grouped by
  reconfiguration sequence length $\ell$, on the opt-available population.
  Column ``\#'' is the number of such instances solved by at least one decomposition.}
  \label{tab:reconflen}
  \begin{tabular}{llrrrr}
    \hline
    variant & length $\ell$ & \# & \OrigDecom & \HeurDecom & \OptDecom \\
    \hline
    \multirow{7}{*}{shortest}
      & $1 \le \ell \le 10$        & $86$  & $83$  & $85$  & $86$  \\
      & $10 < \ell \le 10^{2}$     & $33$  & $28$  & $32$  & $33$  \\
      & $10^{2} < \ell \le 10^{3}$ & $47$  & $47$  & $47$  & $47$  \\
      & $10^{3} < \ell \le 10^{4}$ & $27$  & $27$  & $27$  & $27$  \\
      & $10^{4} < \ell \le 10^{5}$ & $13$  & $13$  & $13$  & $13$  \\
      & $\ell > 10^{5}$            & $2$   & $2$   & $2$   & $2$   \\
      \cline{2-6}
      & total                      & $208$ & $200$ & $206$ & $208$ \\
    \hline
    \multirow{7}{*}{farthest}
      & $1 \le \ell \le 10$        & $63$  & $62$  & $63$  & $63$  \\
      & $10 < \ell \le 10^{2}$     & $66$  & $57$  & $63$  & $66$  \\
      & $10^{2} < \ell \le 10^{3}$ & $43$  & $43$  & $43$  & $43$  \\
      & $10^{3} < \ell \le 10^{4}$ & $31$  & $31$  & $31$  & $31$  \\
      & $10^{4} < \ell \le 10^{5}$ & $12$  & $12$  & $12$  & $12$  \\
      & $\ell > 10^{5}$            & $3$   & $3$   & $3$   & $3$   \\
      \cline{2-6}
      & total                      & $218$ & $208$ & $215$ & $218$ \\
    \hline
  \end{tabular}
\end{table}

The longest sequences are reached precisely because they occur in narrow graphs.
Figure~\ref{fig:reconf-length} plots, for every instance solved under
\HeurDecom, the end-to-end solving time against $\ell$, with each point colored
by the width of its decomposition.
Two regimes are visible.
The long-sequence instances, forming the diagonal band on the right up to
$\ell\approx2.2\times10^{5}$, all have small width and are solved cheaply,
with time growing only gently with $\ell$.
The instances that approach the $30$-minute limit, by contrast, have short
sequences but large width.
Difficulty is therefore governed by width, not by sequence length: once the
width is controlled for, solving time grows roughly linearly in $\ell$ (the
fitted exponent is about $0.9$), so the per-step cost, not the number of steps,
is what makes an instance hard.
This is why \textsc{ddreconf} reaches very long sequences: they live in narrow
graphs where each transition step is inexpensive.
The most extreme case is a power-network instance with $199$ vertices, $302$
edges, and width $5$, whose shortest variant admits a sequence of length
$\ell=221{,}003$ and is solved in about $1{,}700$ seconds (its farthest variant,
$\ell=221{,}016$, is comparable).
A solver that increases the target length one step at a time would require on the
order of $2\times10^{5}$ iterations to certify such a sequence.

\begin{figure}[t]
  \centering
  \includegraphics[width=\linewidth]{reconf_length_time.pdf}
  \caption{End-to-end solving time versus reconfiguration sequence length $\ell$
  for instances solved within $30$ minutes under \HeurDecom, for the
  (a) shortest and (b) farthest variant. Each point is one instance, colored by
  the width of its decomposition. Solving time stays low for the longest
  sequences, which all have small width, while the instances near the time limit
  have short sequences but large width: width, not sequence length, governs
  difficulty.}
  \label{fig:reconf-length}
\end{figure}





\begin{comment}
\section{Analysis}
\label{sec:analysis}

\subsection{Bottlenecks of ZDD-Based Solvers}
\begin{itemize}
\item Solution-space ZDD construction tends to produce large numbers of ZDD nodes.
\item ZDD node counts also grow during transition computation in reconfiguration search.
\item Include plots illustrating these bottlenecks if available.
\end{itemize}

\subsection{Effect on Solution-Space ZDD Size}
\begin{itemize}
\item Analyze how path decomposition reduces the number of ZDD nodes at construction time.
\item Compare ZDD sizes with and without path-decomposition-based preprocessing across benchmark instances.
\end{itemize}

\subsection{Effect on ZDD Size During Reconfiguration}
\begin{itemize}
\item Analyze how path decomposition affects intermediate ZDD sizes during transition computation.
\item Explain why smaller intermediate ZDDs lead to faster reconfiguration search.
\end{itemize}

\section{Experimental Results}
\label{sec:experiments}

\subsection{Comparison with \textsc{ddreconf}}
\begin{itemize}
\item Experimental setup: hardware, benchmark suites, time limits, and parameter settings.
\item Compare runtime, memory usage, and solved instance counts between the proposed solver and \textsc{ddreconf}.
\item Summarize the overall performance improvement.
\end{itemize}

\subsection{Newly Solved Instances}
\begin{itemize}
\item Table listing benchmark instances from the CoRe Challenge 2022 and 2023 solved for the first time by our solver.
\item Brief discussion of the characteristics of these instances.
\item \KJ{TODO}{Replace ``xxx'' with the actual number of newly solved instances.}
\end{itemize}
\end{comment}

\section{Conclusion}
\label{sec:conclusion}
\begin{itemize}
\item Summarize the finding that path decomposition is an effective preprocessing technique for ZDD-based ISR solvers.
\item Emphasize that heuristic path decompositions are sufficient for practical performance gains.
\item Mention the newly solved CoRe Challenge instances.
\item Future work: extension to other reconfiguration problems, tighter integration with other ZDD heuristics, etc.
\end{itemize}

\bibliographystyle{IEEEtran}
\bibliography{bib}

\end{document}
